% from ega1-4-fr.tex
\oldpage[I]{121}préschéma (resp. sous-préschéma fermé) de $X$
s'identifie canoniquement à un sous-préschéma (resp. sous-préschéma fermé) de
$X$.
\end{proposition}

Une partie localement fermée d'un sous-espace localement fermé de $X$ étant
un sous-espace localement fermé de $X$, il est clair
(\hyperref[I.4.1.5-fr]{4.1.5}) que la question est locale et qu'on peut donc
supposer $X$ affine~; la proposition résulte alors de l'identification
canonique de $A/\mathfrak{J}'$ et de
$(A/\mathfrak{J})/(\mathfrak{J}'/\mathfrak{J})$ lorsque
$\mathfrak{J}$, $\mathfrak{J}'$ sont deux idéaux d'un anneau $A$ tels que
$\mathfrak{J}\subset\mathfrak{J}'$.

Nous ferons toujours par la suite l'identification précédente.

\begin{env}[4.1.7]
\label{I.4.1.7-fr}
Soit $Y$ un sous-préschéma d'un préschéma $X$, et désignons par $\psi$
l'injection canonique $Y\to X$ des \emph{espaces sous-jacents}~; on sait que
l'image réciproque $\psi^*(\mathscr{O}_X)$ est la restriction
$\mathscr{O}_X|Y$ (\hyperref[0.3.7.1-fr]{0, 3.7.1}). Si, pour tout $y\in Y$,
on désigne par $\omega_y$ l'homomorphisme canonique
$(\mathscr{O}_X)_y\to(\mathscr{O}_Y)_y$, ces homomorphismes sont les
restrictions aux fibres d'un homomorphisme \emph{surjectif} $\omega$ de
faisceaux d'anneaux $\mathscr{O}_X|Y\to\mathscr{O}_Y$~: il suffit en effet de
le vérifier localement sur $Y$, c'est-à-dire que l'on peut supposer $X$
affine et le sous-préschéma $Y$ fermé~; si dans ce cas $\mathscr{I}$ est le
faisceau d'idéaux dans $\mathscr{O}_X$ qui définit $Y$, les $\omega_y$ ne sont
autres que les restrictions aux fibres de l'homomorphisme
$\mathscr{O}_X|Y\to(\mathscr{O}_X/\mathscr{I})|Y$.

Nous avons donc défini un \emph{monomorphisme d'espaces annelés}
(\hyperref[0.4.1.1-fr]{0, 4.1.1}) $j=(\psi,\omega^b)$ qui est évidemment un
morphisme $Y\to X$ de préschémas (\hyperref[I.2.2.1-fr]{2.2.1}), et que nous
appellerons le \emph{morphisme d'injection canonique}.

Si $f:X\to Z$ est un morphisme, nous dirons encore que le morphisme composé
$Y\xrightarrow{j}X\xrightarrow{f}Z$ est la \emph{restriction} de $f$ au
sous-préschéma $Y$.
\end{env}

\begin{env}[4.1.8]
\label{I.4.1.8-fr}
Conformément aux définitions générales (T, I, 1.1), nous dirons qu'un
morphisme de préschémas $f:Z\to X$ est \emph{majoré} par le morphisme
d'injection $j:Y\to X$ d'un sous-préschéma $Y$ de $X$ si $f$ se factorise en
$Z\xrightarrow{g}Y\xrightarrow{j}X$, où $g$ est un morphisme de préschémas~;
$g$ est nécessairement \emph{unique} puisque $j$ est un monomorphisme.
\end{env}

\begin{proposition}[4.1.9]
\label{I.4.1.9-fr}
Pour qu'un morphisme $f:Z\to X$ soit \emph{majoré par un morphisme
d'injection} $j:Y\to X$, il faut et il suffit que $f(Z)\subset Y$ et que,
pour tout $z\in Z$, si on pose $y=f(z)$, l'homomorphisme
$(\mathscr{O}_X)_y\to\mathscr{O}_z$ correspondant à $f$ se factorise en
$(\mathscr{O}_X)_y\to(\mathscr{O}_Y)_y\to\mathscr{O}_z$ (ou, ce qui revient
au même, que le noyau de $(\mathscr{O}_X)_y\to\mathscr{O}_z$ contienne celui
de $(\mathscr{O}_X)_y\to(\mathscr{O}_Y)_y$).
\end{proposition}

Les conditions sont évidemment nécessaires. Pour voir qu'elles sont
suffisantes, on peut se ramener au cas où $Y$ est un sous-préschéma
\emph{fermé} de $X$, en remplaçant au besoin $X$ par un ouvert $U$ tel que
$Y$ soit fermé dans $U$ (\hyperref[I.4.1.3-fr]{4.1.3})~; $Y$ est alors défini
par un faisceau quasi-cohérent $\mathscr{I}$ d'idéaux de $\mathscr{O}_X$.
Posons $f=(\psi,\theta)$, et soit $\mathscr{J}$ le faisceau d'idéaux de
$\psi^*(\mathscr{O}_X)$, noyau de
$\theta^\sharp:\psi^*(\mathscr{O}_X)\to\mathscr{O}_Z$~; compte tenu des
propriétés du foncteur $\psi^*$ (\hyperref[0.3.7.2-fr]{0, 3.7.2}),
l'hypothèse entraîne que pour tout $z\in Z$, on a
$(\psi^*(\mathscr{I}))_z\subset\mathscr{J}_z$, et par suite
$\psi^*(\mathscr{I})\subset\mathscr{J}$. Par suite, $\theta^\sharp$ se
factorise en
\[
  \psi^*(\mathscr{O}_X)\longrightarrow
  \psi^*(\mathscr{O}_X)/\psi^*(\mathscr{I})
  =\psi^*(\mathscr{O}_X/\mathscr{I})
  \xrightarrow{\omega}\mathscr{O}_Z,
\]
la première flèche étant l'homomorphisme canonique. Soit $\psi'$ l'application
continue $Z\to Y$ coïncidant avec $\psi$~; il est clair que l'on a
${\psi'}^*(\mathscr{O}_Y)=\psi^*(\mathscr{O}_X/\mathscr{I})$~; d'autre part,
$\omega$ est évidemment un homomorphisme local, donc
$g=(\psi',\omega^b)$ est un morphisme $Z\to Y$
